\begin{figure}[!bp]
\centering
\begin{tikzpicture}[scale=0.85,every node/.style={font=\small}]
 \draw[thick] (0,0)--node[above] {$\ell$}(2,0);
 \fill (0,0) circle (2pt) node[left] {$u$};
 \fill (2,0) circle (2pt) node[right] {$v$};
 \draw[-{Stealth},thick] (2.8,0)--(4,0);
 \coordinate (u) at (5,0); \coordinate (v) at (9,0);
 \foreach \h/\lab in {1.4/1,0.45/2,-1.4/j}{
   \draw[thick] (u)--(7,\h)--(v);
   \fill (7,\h) circle (2pt) node[above] {$w_{\lab}$};
 }
 \node at (5.75,1.0) {$\ell/2$};
 \node at (8.25,1.0) {$\ell/2$};
 \node at (7,-0.5) {$\vdots$};
 \fill (u) circle (2pt) node[left] {$u$};
 \fill (v) circle (2pt) node[right] {$v$};
 \node at (1,-1.9) {one old edge};
 \node at (7,-1.9) {countably many two-edge paths};
\end{tikzpicture}
\caption{One normalized diamond replacement, with only three of the
countably many paths shown. Each new edge has length $\ell/2$.
The old endpoints remain at distance $\ell$, and any two distinct new
middle vertices also have distance $\ell$. The drawing is schematic.}
\label{fig:diamond}
\end{figure}
